% =====================================================================
%  07-trabajo-futuro.tex — Trabajo a futuro
% =====================================================================
\chapter{Trabajo a futuro}\label{cap:trabajo-futuro}

Las líneas siguientes parten de las limitaciones de la
\cref{sec:limitaciones} y de los resultados del Capítulo~\ref{cap:pruebas}.

\begin{itemize}
  \item \textbf{Agilidad criptográfica y firma poscuántica.} Evaluar los
        esquemas de firma poscuántica estandarizados por el NIST, solos o en
        modo híbrido con RSA, y medir su efecto en el tamaño del paquete de
        evidencia, en el tiempo de firma en el navegador y en la verificación
        independiente. El formato del paquete debe admitir el cambio de
        algoritmo sin invalidar las firmas ya emitidas.
  \item \textbf{Acreditación de la autoridad.} Estudiar los requisitos
        técnicos, organizativos y jurídicos para que la autoridad certificadora
        y la de sellado operen como Prestador de Servicios de Certificación
        acreditado, o se integren con uno mediante el adaptador previsto, y
        emitan constancias conforme a la NOM-151.
  \item \textbf{Identidad del firmante.} Investigar la vinculación del
        certificado interno con una identidad acreditada por un tercero, como
        la e.firma, para que la firma atribuya el acto a una persona y no solo
        a una cuenta del sistema.
  \item \textbf{Custodia de claves.} Comparar la custodia local en el
        navegador con módulos de seguridad en hardware y con esquemas de firma
        con umbral, incluida la recuperación de una clave perdida sin que el
        servicio la conozca.
  \item \textbf{Transparencia de la bitácora.} Anclar la cabeza de la cadena
        de auditoría en un registro externo y estudiar registros de
        transparencia con pruebas de inclusión y de consistencia que un
        tercero pueda auditar.
  \item \textbf{Reglas procesales.} Formalizar el cómputo de plazos como un
        modelo verificable, ampliar el corpus con revisión jurídica calificada
        y extenderlo a otras materias y fueros.
  \item \textbf{Escala y operación.} Caracterizar el comportamiento bajo
        carga y el aislamiento entre varios despachos en una misma
        instalación.
\end{itemize}
